\chapter{Autres expériences significatives}

Au-delà de mon parcours professionnel principal, mon engagement dans l'innovation pédagogique numérique se manifeste par une participation active à des projets collaboratifs d'envergure nationale. Ces expériences complémentaires révèlent une démarche de recherche-action et d'ingénierie de formation qui s'inscrit parfaitement dans les exigences du Master MEEF PIF.

\section{Contributeur actif à la Forge des Communs Numériques Éducatifs}

\subsection{Contexte et présentation de l'initiative}

La Forge des Communs Numériques Éducatifs (LaForgeEdu) est une plateforme collaborative nationale hébergée sur GitLab, dédiée au développement et au partage de ressources pédagogiques sous licence libre \parencite{forge_educative}. Cette initiative s'inscrit dans la Stratégie du numérique pour l'éducation 2023-2027 du ministère de l'Éducation nationale.

\textbf{Mon rôle :} Contributeur développeur et accompagnateur pédagogique depuis 2022.

\textbf{Mission principale :} Développement d'applications éducatives libres et accompagnement d'enseignants dans la création de leurs propres outils numériques.

Mes développements sur la Forge ne sont pas de simples utilitaires, mais s'apparentent à la conception d'Environnements Informatiques pour l'Apprentissage Humain (EIAH). Par exemple, le "Générateur d'Exercices Python" intègre des fonctionnalités de diagnostic cognitif et de rétroaction (feedback) élaborée. En m'appuyant sur les traces d'apprentissage (learning analytics) laissées par les élèves lors de leurs interactions avec l'outil, je peux ajuster mes interventions didactiques. Cette approche rejoint les problématiques traitées lors des conférences EIAH ou TICE, notamment sur la capacité des artefacts numériques à soutenir l'autorégulation de l'apprenant dans sa zone proximale de développement.

\subsection{Activités et réalisations concrètes}

\textbf{Portfolio de projets développés (données extraites de ebep.educ-ai.fr) :}

\begin{center}
\footnotesize
\begin{tabular}{|p{3cm}|p{5.5cm}|p{3.5cm}|p{4cm}|}
\hline
\textbf{Projet} & \textbf{Description} & \textbf{Technologies} & \textbf{Impact pédagogique} \\
\hline
Générateur d'Exercices Python avec IA & Application web pour générer des exercices de programmation personnalisés & Python, Streamlit, IA & Personnalisation des apprentissages NSI \\
\hline
PYTEUR OS & Plateforme éducative collaborative avec génération d'exercices assistée par IA & Django, IA, PostgreSQL & Gestion pédagogique complète \\
\hline
L'Établi & Application web pour simplifier la gestion de dépôts sur la forge éducative & Vue.js, API GitLab & Démocratisation de l'accès à la forge \\
\hline
Chat avec IA via RAG & Application de chat basée sur la technique RAG pour interroger des documents PDF & IA, RAG, LLM & Recherche documentaire assistée \\
\hline
Pokédex Interactif & Application Python ludique pour l'apprentissage de la programmation en Terminale NSI & Python, Streamlit, SQLite & Gamification de l'apprentissage \\
\hline
\end{tabular}
\end{center}



\section{Recherche-action sur l'Intelligence Artificielle dans l'Éducation}

\subsection{Contexte et enjeux}

\textbf{Groupe de travail spécialisé :} Participation au salon « IA LaForgeEdu » sur Tchap, communauté de « professeurs bidouilleurs d'IA » qui documentent et partagent leurs expérimentations sur l'intégration de l'intelligence artificielle dans l'enseignement.

\textbf{Problématique centrale :} Comment intégrer de manière éthique et pédagogiquement pertinente l'intelligence artificielle dans les pratiques d'enseignement, particulièrement pour l'inclusion des Élèves à Besoins Éducatifs Particuliers (EBEP) et la différenciation pédagogique ?

\subsection{Contributions et réalisations}

\textbf{Plateforme ebep.educ-ai.fr :} Création et animation d'un site de référence dédié à l'intégration de l'IA pour l'accompagnement des Élèves à Besoins Éducatifs Particuliers (EBEP), proposant des outils d'adaptation, d'accessibilité et de différenciation pédagogique.

\begin{center}
\footnotesize
\begin{tabular}{|p{4cm}|p{5cm}|p{7cm}|}
\hline
\textbf{Domaine d'expérimentation} & \textbf{Outils développés / Ressources} & \textbf{Résultats et impact pour les EBEP} \\
\hline
Adaptation et Accessibilité DYS & Générateur d'exercices adaptés, reformulations simplifiées par IA & Différenciation pédagogique, réduction de la charge cognitive, meilleure lisibilité \\
\hline
Assistance documentaire \& FALC & Chat RAG pour simplification en Français Facile à Lire et à Comprendre & Compréhension facilitée des consignes, autonomie accrue des élèves à besoins particuliers \\
\hline
Plateforme collaborative \& EBEP & PYTEUR OS / ebep.educ-ai.fr avec IA intégrée pour la différenciation & Suivi personnalisé des élèves EBEP, création accélérée de supports adaptés \\
\hline
Éthique, sécurité \& inclusion & Réflexion sur l'usage responsable de l'IA pour l'inclusion numérique & Sensibilisation de la communauté éducative, prévention des dérives, équité d'accès \\
\hline
\end{tabular}
\end{center}
